\section{Research Gap and Objectives}

Existing literature provides robust methodologies for analyzing generalization failure; however, it lacks predictive tools applicable prior to deployment.

The main problem is not a lack of awareness, since the issue is well covered in existing studies. Instead, the problem is the lack of a combined process that includes measuring data differences, analyzing individual features, and providing deployment advice.

\textbf{Research Question:}\\
How well can we predict before deployment if clinical machine learning models will perform worse on new patient groups by checking data differences and feature reliability? Also, what cutoff points can be set to decide if it is safe to deploy the model?

\textbf{Objectives:}

\begin{enumerate}
    \item \textbf{Objective 1:} Apply MMD \cite{gretton2012kernel}, KL divergence, and Wasserstein distance methods \cite{duchi2018learning, rahimian2019dro} to clinical data from different settings where performance drops are known. Compare each score with the real loss in balanced accuracy.
    
    \item \textbf{Objective 2:} Use leave-one-context-out permutation importance \cite{whalen2022pitfalls} to mark each predictor as either stable or sensitive to changes. Check these labels against actual performance across different settings, aiming for at least 70\% agreement..
    
    \item \textbf{Objective 3:} Combine Objectives 1 and 2 into a decision flowchart with four options: deploy with normal monitoring, deploy with extra monitoring, retrain using stable features, or do not deploy. Test this on at least one separate setting and clearly report its limitations.
\end{enumerate}